\chapter{Conclusiones y trabajos futuros}\label{cap:conclusiones}

Este capítulo de cierre no introduce ninguna cifra nueva: toda afirmación que sigue
remite a una tabla o figura ya presentada y discutida en el capítulo
\ref{cap:resultados}, del que este capítulo es una síntesis argumental y no una nueva
fuente de evidencia.

\section{Conclusiones}\label{sec:conclusiones}

El objetivo general —evaluar si una arquitectura híbrida residual LSTM-XGBoost mejora la
predicción de la demanda diaria agregada de transporte público en Madrid frente a
baselines estadísticos y a modelos de una sola etapa— se ha cumplido: la evaluación se
realizó y su respuesta es negativa para la arquitectura. El cumplimiento de cada objetivo
específico de la sección \ref{sec:objetivos} se acredita por separado:

\begin{itemize}
  \item \textbf{OE1, ingesta y matriz anti-fuga.} \emph{Qué se hizo:} tres fuentes crudas
    unificadas en 1.310 días continuos y una matriz de 161 variables predictoras con
    guardias que fallan de forma explícita (capítulo \ref{cap:datos_features}).
    \emph{Resultado:} ninguna guardia ha detectado una fuga en la matriz publicada; las
    dos que durante varias fases eran inalcanzables se corrigieron y disparan hoy bajo
    prueba (Anexo \ref{cap:anexo2}). \emph{Cumplido.} \emph{Limitación:} la ausencia
    total de fuga no puede afirmarse sin una auditoría exhaustiva externa (sección
    \ref{sec:limitaciones}).
  \item \textbf{OE2, baselines y SARIMAX.} \emph{Qué se hizo:} cuatro baselines ingenuos y
    un SARIMAX(2,1,2)$\times$(0,1,2,7) seleccionado por AIC de entrenamiento, todos a un
    día vista (sección \ref{sec:baselines}). \emph{Resultado:} SARIMAX alcanza un MAE de
    test de 163.627 y es uno de los dos modelos de una sola etapa más precisos: se
    incorpora como un modelo de referencia competitivo y metodológicamente relevante
    (PI1, PI2). \emph{Cumplido.} \emph{Limitación:}
    el orden elegido no está fuertemente identificado (cinco órdenes en 0,8 unidades de
    AIC).
  \item \textbf{OE3, etapa 1 con predicciones OOF.} \emph{Qué se hizo:} LSTM univariante
    con esquema \emph{walk-forward} expansivo de cinco folds y escalador reajustado por fold
    (sección \ref{sec:esquema_oof}). \emph{Resultado:} 689 predicciones OOF sobre 889 filas;
    el fold 1 resultó degenerado y se excluyó tras comprobar su cobertura.
    \emph{Cumplido.} \emph{Limitación:} el residuo de entrenamiento procede de modelos con
    337-748 filas mientras que el de inferencia procede del modelo con 889, un
    desplazamiento inherente al \emph{stacking} OOF.
  \item \textbf{OE4, etapa 2 y controles.} \emph{Qué se hizo:} XGBoost sobre 552 residuos
    OOF con las 161 variables, comparado con la etapa 1 sola y con un XGBoost directo sobre
    la misma matriz (secciones \ref{sec:etapa2_xgboost} y \ref{sec:modelos_control}).
    \emph{Resultado:} la etapa 2 reduce el MAE de su etapa 1 en un 16,6\,\% en test, pero
    el híbrido queda por detrás de SARIMAX y de XGBoost independiente en ambas
    particiones (PI3, PI4). \emph{Cumplido en su formulación} —el contraste se realizó y
    los controles aislaron la causa—, \emph{con resultado negativo}. \emph{Limitación:}
    5 festivos y 3 puentes en test impiden concluir sobre los días irregulares.
  \item \textbf{OE5, ensamble, sensibilidad e interpretación.} \emph{Qué se hizo:} dos
    experimentos acotados y una lectura a la luz de \textcite{granger1989combining}
    (sección \ref{sec:analisis_sensibilidad}). \emph{Resultado:} el repesado no mejora los
    días que perseguía; el ensamble aritmético de SARIMAX y XGBoost independiente es el
    modelo más preciso del estudio (PI5). \emph{Cumplido.} \emph{Limitación:} la
    preferencia entre las dos variantes del ensamble se formuló tras conocer sus cifras de
    test (sección \ref{sec:limitaciones}).
\end{itemize}

El hallazgo empírico central aparece en la evaluación de la corrección residual:
\textbf{el híbrido, tal como queda diseñado, no supera a SARIMAX ni a
XGBoost independiente en ninguna de las dos particiones} (PI1 y PI2, tabla
\ref{tab:delta_hibrido}), y tampoco reduce de forma diferencial el error en los días de
calendario irregular que motivaron su diseño (PI3, tabla \ref{tab:desglose_dia_test}).
No es un fallo de ejecución sino un resultado consistente con la teoría de combinación de
pronósticos: \textcite{granger1989combining} exige que los componentes sean
individualmente subóptimos \emph{y} se apoyen en conjuntos de información distintos, y
esta arquitectura no satisface plenamente la segunda condición —los conjuntos de información de sus dos
etapas están anidados, no son complementarios: la etapa 2 añade calendario y meteorología,
pero el pasado reciente que ve la etapa 1 ya está representado en su matriz— como se
argumenta en la sección
\ref{sec:resultado_negativo_teoria}. Las competiciones M4 y M5
\parencite{makridakis2018m4,makridakis2022m5} documentan que el patrón —métodos simples y
combinaciones bien fundamentadas superando a arquitecturas más sofisticadas, sobre todo en
muestras moderadas como las 1.310 observaciones de este proyecto— es esperable y no una
anomalía de esta implementación.

Tampoco debe leerse como una arquitectura sin mérito: PI4 muestra que la etapa 2 cumple
su función, reduciendo el MAE de su propia etapa 1 en un 16,6\,\% en test y un 58\,\% en
festivos; su variable de mayor ganancia, \texttt{feat\_day\_type\_festivo}, es coherente
con que esa corrección proceda de la señal de calendario que una LSTM univariante no puede
ver por construcción (sección \ref{sec:importancia_variables}). El mecanismo que justifica
la hibridación se observa de forma medible; lo insuficiente es la etapa 1 de la que parte.

El modelo que mejor resuelve el problema no requiere ninguna red neuronal: PI5 muestra
que un ensamble aritmético de SARIMAX y XGBoost independiente —sin reentrenamiento, sin
etapa LSTM y sin ningún paso de ajuste— alcanza un MAE de test de 139.725 en su variante
equiponderada (el ensamble 50/50) y de 139.681 en la ponderada por el inverso del MAE
(el ensamble inverso al MAE), en torno a un 10,5\,\% por debajo del mejor de sus
componentes. Es coherente con la identidad de
\textcite{bates1969combination} aplicada a la correlación de 0,576 entre los residuos de ambos
sobre test (sección \ref{sec:mecanismo_ensemble}): inferior a la unidad porque fallan en
subconjuntos de días sistemáticamente distintos —el primero en laborables ordinarios, el
segundo en fin de semana y festivos—, que es lo que hace ventajoso promediarlos.

La lectura conjunta de ambos resultados —el híbrido secuencial que no mejora porque sus
etapas observan conjuntos de información anidados, y el ensamble paralelo que sí mejora, en
consonancia con la combinación de información parcialmente descorrelacionada— es la
aportación empírica central de este TFM:
una auditoría con protocolo \emph{out-of-fold} estricto en toda su cadena de las
condiciones exactas bajo las cuales cada estrategia de combinación funciona y bajo cuáles
no, sobre una serie real de demanda de transporte público agregada y multimodal.

\section{Limitaciones}\label{sec:limitaciones}

\textbf{Agregación a nivel de red, sin corte espacial.} Este proyecto modela la demanda
\texttt{total} del sistema CRTM como una única serie diaria, con el desglose por operador
solo como variable de apoyo (rezagos, EDA) y no como dimensión espacial del modelo. A
diferencia de \textcite{cardozo2012spatial}, que modelan las entradas estación
por estación, esta memoria no ofrece resolución espacial alguna: de sus resultados no cabe
concluir nada sobre una línea, una estación o un corredor concreto. La elección fue
deliberada (sección \ref{sec:movilidad_urbana}), pero es una limitación de alcance.

\textbf{Tamaño muestral reducido en los días de calendario más interesantes.} Con 1.310
observaciones diarias en total, la muestra es modesta para un modelo LSTM, y se vuelve
particularmente escasa precisamente en los subgrupos de calendario cuyo comportamiento
resulta más interesante para este estudio: la partición de test contiene solo 5 días
festivos y 3 días de laborable con puente (tabla \ref{tab:desglose_dia_test}). Todas
las afirmaciones de esta memoria sobre el comportamiento de los modelos en esos dos
grupos, incluida la mejora del 58\,\% de la etapa 2 sobre su etapa 1 en festivos, se han
etiquetado consistentemente como indicativas y no concluyentes; la base más sólida para
generalizar sobre el comportamiento estructural frente al calendario sigue siendo el
diagnóstico exploratorio de la fase 4 sobre el periodo completo (48 festivos y 54
puentes; la detección de calendario identifica 55 días puente sobre los 1.310 días
totales, de los que 54 caen fuera de los 28 días de calentamiento excluidos del
entrenamiento), no las particiones reducidas de evaluación final.

\textbf{Recorte de 28 días por calentamiento de las variables de retardo.} El retardo y la
media móvil de mayor ventana (28 días) obligan a descartar los 28 primeros días de la
serie de entrenamiento (2023-01-01 a 2023-01-28) antes de que la matriz de
características quede libre de valores nulos (sección \ref{sec:particiones_fuga}). Esta
política, aplicada únicamente sobre el bloque de entrenamiento y verificada para no
alcanzar nunca las particiones de validación o test, reduce el conjunto de entrenamiento
efectivo de 917 a 889 días: un coste aceptado y documentado, pero un coste real sobre
el volumen de datos disponible para el ajuste de los modelos.

\textbf{Naturaleza provisional de los registros de demanda del CRTM.} El propio fichero
fuente de demanda (\texttt{CRTM\_Evolucion\_demanda\_diaria.xlsx}) advierte
explícitamente que sus cifras son provisionales y sujetas a revisión posterior por el
consorcio. Este TFM no tiene control sobre esa advertencia ni forma de cuantificar la
magnitud de una eventual revisión futura; toda cifra de demanda y todo error de modelo
reportado en esta memoria hereda, por tanto, la incertidumbre de la fuente original.

\textbf{Exclusión deliberada de la meteorología del mismo día.} Se excluye por diseño la
meteorología concurrente del día predicho, empleando solo sus rezagos de 1, 2, 3 y 7 días
y su media móvil de 7 (sección \ref{sec:promocion_meteo}), para no asumir en evaluación un
pronóstico atmosférico perfecto que no existe en producción. La decisión es conservadora y
correcta para el objetivo del trabajo, pero implica que ninguna cifra de esta memoria
estima la mejora alcanzable con un pronóstico meteorológico operativo de calidad: esa cota
queda fuera de alcance y se retoma como línea futura.

\textbf{Veredictos parcialmente ambiguos en la anatomía del residuo.} El diagnóstico del
residuo de la etapa 1 (sección \ref{sec:anatomia_residuo_etapa1}) evaluó cuatro criterios
de aceptación fijados por adelantado, y solo dos cayeron con claridad en uno de los dos
extremos: la concentración del error en el calendario irregular (``la estructura
persiste'', factor 6,52 en festivos) y la ausencia de señal meteorológica retardada
detectable (``ruido cuasi-blanco'', ninguna correlación parcial significativa tras
Benjamini-Hochberg). Los otros dos, la memoria temporal general y la estacionalidad
semanal, quedaron en la banda intermedia: la mayor autocorrelación del residuo es 0,166,
por debajo del umbral de 0,20 fijado como evidencia de estructura, pero por encima del de
0,10 que indicaría ruido, y el pico espectral de periodo 7 días es débil por fuga
espectral. En consecuencia, la afirmación de que la etapa 1 deja sin explicar estructura
\emph{semanal} descansa en el tamaño del efecto del perfil por día de la semana
(dispersión de 0,60 desviaciones típicas) y no en los criterios de autocorrelación o
espectrales, que son individualmente inconcluyentes. El hallazgo central de esa sección
(que el residuo retiene estructura de \emph{calendario}, no de clima) no se ve afectado,
pero la caracterización de su componente temporal fino es más matizada de lo que un
veredicto binario sugeriría.

\textbf{Veredicto ambiguo en el control de paridad informativa con matriz completa.} El
control de paridad informativa (sección \ref{sec:paridad_lstm}) evaluó dos ámbitos de
variables predictoras con criterios de cierre de hueco fijados por adelantado. El ámbito
\texttt{calendario} ($k=17$) cayó con claridad en \emph{limitación arquitectónica}: cierra
solo un 16,2\,\% del hueco hasta XGBoost independiente, con la banda mín-máx de cinco
semillas por completo dentro de la zona de ese veredicto. Pero el ámbito \texttt{completo}
($k=161$) cerró un 31,0\,\% y su banda (354.235--391.630) cruza el umbral del 25\,\%, de
modo que su veredicto se degrada a \emph{mixto/ambiguo} por la regla de cruce de banda
(objeción O6). La conclusión de la sección —que la etapa 1 pierde por ser recurrente sobre
esta serie y a este tamaño de muestra, y no por ser univariante— se apoya en el ámbito
\texttt{calendario}, que es el más nítido para la afirmación de ceguera al calendario, y en
la asimetría de interpretación O4 (la LSTM en paridad ve estrictamente más información que
XGBoost, así que un cierre parcial es evidencia débil). Aun así, no puede descartarse con
la matriz completa que una fracción del hueco se explique por acceso a la información y no
solo por la arquitectura; discriminar ambas causas con más semillas o un diseño que iguale
también la representación de entrada queda fuera del alcance de este TFM.

\textbf{Elección entre las dos variantes del ensamble tras conocer sus cifras de test.}
No existía una regla previa para preferir el ensamble 50/50 al
ensamble inverso al MAE: ambas se calcularon en la misma ejecución y la
preferencia se formuló con las dos cifras de test a la vista (sección
\ref{sec:por_que_ensemble}). El test no interviene en ningún peso y la elección no lo
optimiza —la equiponderada es la peor de las dos tanto en validación como en test—, y
elegir entre ambas por su MAE de validación sería circular en un sentido preciso: no por
ajustar pesos en validación y evaluar en test, que es el procedimiento estándar, sino porque
los pesos de una de las dos variantes se ajustaron sobre esa misma partición y su cifra de
validación es, por ello, levemente optimista. Aun así, la decisión de presentación se tomó a
posteriori y se registra como tal; la diferencia entre ambas, 44 viajes diarios de MAE,
carece de relevancia práctica.

\textbf{Alcance del \emph{bootstrap} por bloques.} La única inferencia estadística de la
memoria (sección \ref{sec:desglose_incertidumbre}) usa un \emph{bootstrap} pareado de
bloques móviles de 7 días sobre 197 días de test. La longitud de bloque se fijó por la
autocorrelación semanal dominante, no se estimó, y no captura una dependencia de escala
mayor (mensual o estacional); con un solo periodo de test, los intervalos describen la
variabilidad de remuestreo dentro de ese periodo, no la de otros periodos, y los subgrupos
solapados son contrastes dependientes. Por eso las diferencias entre los modelos punteros
de la tabla maestra se presentan sin contraste de significación.

\textbf{Imposibilidad de asegurar la ausencia total de fuga sin una auditoría
exhaustiva.} Los controles anti-fuga cubren las formas de fuga identificadas en el diseño
(retardos y ventanas desplazadas, particiones cronológicas, escalado por fold,
predicciones OOF, exclusión de operadores y meteorología del mismo día) y cada uno tiene
una prueba negativa que lo obliga a disparar. Que ninguno haya detectado una fuga en la
matriz publicada no equivale a demostrar que no exista ninguna: un tipo de fuga no
previsto en el diseño no tendría guardia, y la incidencia de las dos guardias
inalcanzables del Anexo \ref{cap:anexo2} muestra que la propia batería puede tener
defectos silenciosos. La afirmación que sostiene esta memoria es la prudente: se han
implementado controles específicos contra las principales formas de fuga, y se han
verificado sobre el código.

\textbf{Dependencia de una única serie y de un único ámbito geográfico.} Todos los
resultados proceden de una serie diaria de 1.310 observaciones del sistema CRTM de la
Comunidad de Madrid en 2023-2026. Ni el resultado negativo del híbrido ni la ventaja del
ensamble se han replicado sobre otra red, otro periodo ni otra escala temporal; su
generalización a otros sistemas o a demanda horaria es una hipótesis que este trabajo no
contrasta.

\section{Trabajos futuros}\label{sec:trabajos_futuros}

Las siguientes líneas de trabajo se enuncian como extensiones metodológicas concretas y
justificadas por los propios resultados de este TFM; ninguna de ellas ha sido ejecutada
ni evaluada en el marco de este proyecto.

\textbf{Modelado desagregado por operador.} Este TFM modela la demanda \texttt{total}
del sistema, preservando \texttt{metro}, \texttt{emt}, \texttt{carretera} y
\texttt{cercanias} únicamente como variables de apoyo. Una extensión natural sería
replicar el mismo pipeline de ingeniería de características y la misma arquitectura de
evaluación (baselines, LSTM \emph{out-of-fold}, XGBoost residual, ensamble) de forma
independiente para cada uno de los cuatro operadores, lo que permitiría dos preguntas
que este TFM no puede responder: si la ventaja del ensamble sobre el híbrido se mantiene
a una escala de demanda menor y con patrones de calendario específicos de cada modo —el
perfil de Cercanías, más ligado a la movilidad laboral pendular, es previsiblemente
distinto del perfil de la red de EMT—, y si un modelo conjunto multivariante que
comparta representación entre operadores aporta algo que cuatro modelos independientes
no capturan.

\textbf{Integración de pronósticos meteorológicos operativos.} Tal como se señala en la
limitación correspondiente, este TFM emplea exclusivamente meteorología retardada para
evitar asumir un pronóstico perfecto. Una extensión natural, y metodológicamente sólida,
sería sustituir esos rezagos por pronósticos meteorológicos
reales de horizonte 1-7 días obtenidos de un proveedor operativo (por ejemplo, la propia
API de pronóstico de Open-Meteo, complementaria a su API histórica ya empleada en este
proyecto), evaluando el sistema bajo la incertidumbre de pronóstico real en lugar de bajo
el valor observado retardado. Esto permitiría cuantificar, con evidencia y no por
extrapolación, la brecha entre el límite superior bajo pronóstico perfecto y el
desempeño alcanzable en un escenario de producción genuino.

\textbf{Arquitecturas neuronales basadas en parches o transformers temporales.} La etapa
1 de este TFM es deliberadamente una LSTM univariante simple, elegida para aislar con
claridad qué puede aportar la dinámica temporal pura antes de introducir información
exógena en la etapa 2. Arquitecturas más recientes de pronóstico de series temporales,
basadas en mecanismos de atención y en la segmentación de la serie en parches temporales,
se han propuesto precisamente para representar dependencias de múltiples escalas —semanal
y anual simultáneamente— sin la memoria secuencial de una red recurrente clásica; si sobre
esta serie lo consiguen es una cuestión empírica abierta. Sustituir la etapa 1 de este TFM
por una arquitectura de este
tipo, entrenada bajo el mismo protocolo estrictamente \emph{out-of-fold} que gobierna
todo este proyecto, permitiría comprobar si una etapa 1 más expresiva cierra la brecha
frente a SARIMAX que este TFM documenta como su hallazgo central, sin renunciar en
ningún momento al rigor de validación que hace de esa comprobación un resultado
comparable y no una nueva fuente de sesgo optimista.

\section{Implicaciones para la planificación}\label{sec:implicaciones_planificacion}

Sin introducir ninguna evidencia nueva, los resultados del capítulo
\ref{cap:resultados} admiten una lectura orientada a la explotación operativa, que es el
destino último que motivó este TFM. El ensamble seleccionado predice la demanda diaria
agregada de la red con un MAPE de test del \textbf{3,13\,\%} (cuadro \ref{tab:comparacion_test}),
un margen compatible con los usos de planificación en los que la unidad de decisión es la
red completa y el horizonte es el día siguiente, el único evaluado en este trabajo:
dimensionamiento de la oferta de material y de
turnos de personal, previsión de carga agregada para la programación de frecuencias, y
detección temprana de desviaciones frente a lo esperado. Para esos usos, la recomendación
que se desprende de este trabajo es doble y deliberadamente austera: emplear la
combinación aritmética de SARIMAX y XGBoost independiente descrita en la sección
\ref{sec:por_que_ensemble}, y no una arquitectura recurrente ---el capítulo
\ref{cap:resultados} muestra que la etapa LSTM no solo no aporta precisión, sino que su
coste de entrenamiento, ajuste y mantenimiento no se traduce en ninguna ventaja
medible---; y fijar los pesos a partes iguales, porque la variante ponderada no mejora de
forma apreciable y sí añade una dependencia de la partición de validación que habría que
rehacer con cada reentrenamiento.

Esa recomendación tiene límites que deben acompañarla siempre, y que son exactamente los
de la sección \ref{sec:limitaciones}. El error no se reparte de forma uniforme: se
concentra en los días de calendario irregular ---festivos y puentes (cuadro
\ref{tab:desglose_dia_test})---, que son precisamente aquellos en los que una decisión de
oferta equivocada resulta más costosa, de modo que las cifras de error agregadas no deben
trasladarse sin más a la planificación de esos días. El modelo no ofrece resolución
espacial alguna por debajo del agregado de red, ni predice la demanda de un operador
concreto, por lo que no sustituye a ningún instrumento de planificación de línea,
corredor o estación. Y, al excluir por diseño la meteorología del mismo día, su desempeño
es el alcanzable sin pronóstico meteorológico: incorporarlo es la vía de mejora
inmediata, y es la segunda de las líneas enunciadas en la sección
\ref{sec:trabajos_futuros}.

\section{Contribución a los Objetivos de Desarrollo Sostenible}\label{sec:ods}

Este trabajo se vincula a la Agenda 2030 de Naciones Unidas a través de dos objetivos.
Conviene precisar el alcance de esa vinculación antes de enunciarla, porque el hallazgo
central de esta memoria es negativo y una contribución declarada sin esa salvedad sería
engañosa: lo que sigue no describe una mejora desplegada, sino una aportación
metodológica.

\textbf{ODS 11, ciudades y comunidades sostenibles (meta 11.2).} El objeto de estudio es
la demanda agregada del sistema de transporte público de Madrid, entre 1,3 y 6,6 millones
de viajes diarios, que es el modo de desplazamiento que la meta 11.2 busca hacer
accesible, asequible y sostenible. La predicción a un día alimenta las decisiones de
dimensionamiento de flota, turnos de personal y programación de frecuencias descritas en
la sección \ref{sec:implicaciones_planificacion}: una previsión más ajustada permite
sostener el nivel de servicio sin sobredimensionar la oferta, que es la forma en que un
sistema de transporte público mejora su acceso sin aumentar su coste. La contribución es
indirecta ---el trabajo no despliega ningún sistema--- y está acotada exactamente por la
sección \ref{sec:limitaciones}: no hay resolución espacial por debajo del agregado de red
y el error se concentra en el calendario irregular, que son los días en los que una
decisión de oferta equivocada resulta más costosa.

\textbf{ODS 13, acción por el clima.} La aportación específica de este trabajo a este
objetivo no está en la precisión alcanzada, sino en su coste. La sección
\ref{sec:paridad_informativa} documenta que la etapa recurrente no aporta precisión
medible, y el mejor resultado del estudio lo obtiene una combinación aritmética de dos
modelos ya entrenados, sin entrenamiento adicional ni aceleración por hardware
especializado. En un campo cuya respuesta por defecto ante un problema de predicción es
escalar la arquitectura, un resultado negativo bien documentado ahorra ese gasto
computacional a quien venga después, y ese ahorro sí es directamente atribuible al
trabajo. No se cuantifica aquí ninguna reducción de emisiones ni se afirma relación
causal alguna entre la calidad de la predicción y la huella ambiental del sistema de
transporte: sostenerlo excedería lo que permiten los datos empleados.

% Ultima pagina del cuerpo (capitulos 1-6). El pie de pagina la usa como total, de modo que
% la extension que cuenta para el limite se lee directamente del folio y no hay que
% recomputarla restando anexos y bibliografia. Debe quedar antes del \appendix de
% 07_anexos.tex, que reinicia la numeracion.
\label{CuerpoLastPage}
